\section{管辖}
\begin{enumerate}
	\item 先定被告再做管辖，被告不知管辖不明，管辖内部先级别后地域。
	\item 步骤。先定被告 \yjt 先中后基 \yjt 不动产 \yjt 复议/限制人身自由 \yjt 被告所在地。
	\item 级别。基层是原则，中院是例外。
	      \begin{enumerate}
		      \item 除非授权，专门法院（海事、铁路、军事法院）一般不审理行政案件。
		      \item \zx{复议维持}，\bs{级别就低}。\zx{原机关决定级别管辖}。\al{被告为乙\xhx{县政府}和\xhx{药监局}，乙县法院对此案有管辖权 \dui}。
	      \end{enumerate}

	      \begin{kk}{中院。重点，每年必考}
		      \begin{enumerate}
			      \item 被告是\zd{县级以上政府和国务院部门}（重中之重每年必考）。xx 部 = 国务院部门。区政府 \ydy 县政府。
			      \item 涉外（含港澳台）。
			      \item 海关、进出口、股票和证券（证交所）。
		      \end{enumerate}
	      \end{kk}
	\item 地域。\zx{原告就被告}。\zd{最初作出行政行为的行政机关所在地}。原告地包括户籍地、经居地、Opt(被限制人身自由地)。
	      \begin{kk}{特别的地域管辖}
		      \begin{enumerate}
			      \item 不动产（因行政行为导致\zd{物权变动}）。不动产所在地（有本按本，没本实际）。\al{房屋登记、土地确权}。
			      \item 便利原告原则。
			            \begin{enumerate}
				            \item \zd{经过复议。原机关 \hz 复议机关所在地}。永远这俩地。
				            \item 限制人身自由。原告 + 被告。\bs{关了的人诉关了的强制措施}。行政拘留是处罚，不是强制措施。
				                  \zd{对限制人身自由的行政强制措施不服}。
				                  \begin{enumerate}
					                  \item \al{甲住 A 县，被 B 县公安关在 C 县。向 D 区复议，维持，甲起诉。ABCD}。
					                  \item \al{戒毒强制拘留 \dui}。
				                  \end{enumerate}
			            \end{enumerate}
			      \item 跨区域管辖最高法批准。
		      \end{enumerate}
	      \end{kk}
	\item \bs{如何确定本案的地域管辖}。\zx{……不属于特殊地域管辖情形……}。
\end{enumerate}
